
\medskip
\phantomsection\label{cor:sga3-VIB-6.5.4}%
\noindent\textbf{Corollary 6.5.4.---}
\emph{Let \(X\) be a smooth \(S\)-group scheme of finite
presentation with connected fibers, let \(Y\) be a group scheme
of finite presentation over \(S\), and let \(i:Y\to X\) be a
monomorphism of \(S\)-group schemes.}
\begin{enumerate}
\item[\textup{(i)}]\phantomsection\label{concl:sga3-VIB-6.5.4-i}
  \emph{The functor
  \[
    \prod_{X/S}Y/X
  \]
  is representable by a closed subscheme of \(S\) of finite
  presentation.}
\item[\textup{(ii)}]\phantomsection\label{concl:sga3-VIB-6.5.4-ii}
  \emph{If, moreover, \(S\) is quasi-compact, then for all
  sufficiently large \(n\),
  \[
    \prod_{X/S}Y/X
      =\prod_{X_n/S}Y_n/X_n,
  \]
  where \(X_n\) denotes the \(n\)-th infinitesimal neighborhood
  of the unit section \(\epsilon:S\to X\), and
  \(Y_n=X_n\cap Y\).}
\end{enumerate}

The proof is essentially that of
\hyperref[cor:sga3-VIB-6.5.3]{6.5.3}.%
\footnote{\label{note:sga3-VIB-65}Editor's note: the original has
been expanded in what follows, and references to EGA~IV have
been added.}
On the one hand, \(i\) is locally of finite presentation
(EGA~IV\(_1\), 1.4.3). On the other hand, the unit sections of
\(Y\) and \(X\) induce bijective immersions
\[
  S\longrightarrow Y_n
  \qquad\text{and}\qquad
  S\longrightarrow X_n,
\]
and hence isomorphisms of \(S_{\mathrm{red}}\) with
\((Y_n)_{\mathrm{red}}\) and \((X_n)_{\mathrm{red}}\).
Consequently, \(i_n\) is quasi-compact, hence of finite type, and
one has the commutative diagram
\begin{figure}[H]
  \centering
  \phantomsection\label{fig:sga3-VIB-6.5.4-reduced-triangle}%
  \begin{tikzcd}[row sep=large,column sep=large]
    (Y_n)_{\mathrm{red}}
      \arrow[r,"\tau"]
      \arrow[dr,"\sigma"'] &
    Y_n \arrow[d,"i_n"'] \\
    {} &
    X_n
  \end{tikzcd}
\end{figure}
where \(\sigma\) and \(\tau\) are closed immersions and \(\tau\)
is surjective. Since \(i_n\) is separated (being a
monomorphism), it follows that \(i_n\) is proper
(EGA~II, 5.4.3). Thus \(i_n\) is a proper monomorphism of finite
presentation, and therefore a closed immersion
(EGA~IV\(_3\), 8.11.5). It follows from
\hyperref[thm:sga3-VIB-6.2.3]{6.2.3} that
\[
  \prod_{X_n/S}Y_n/X_n
\]
is representable by a closed subscheme \(T_n\) of \(S\), of
finite presentation over \(S\). It remains to prove the last
assertion of \hyperref[cor:sga3-VIB-6.5.4]{6.5.4} when \(S\) is assumed affine as well.

We again reduce immediately to the case where \(S\) is
Noetherian. With the notation of the proof of \hyperref[cor:sga3-VIB-6.5.3]{6.5.3}, it is enough
to prove \(R=T\), or equivalently that
\[
  Y\supset X_n\quad\text{for every }n
  \qquad\Longrightarrow\qquad
  Y=X.
\]
The hypothesis implies that \(i:Y\to X\) is étale at the points
of the unit section of \(Y\) over \(S\). Thus \(Y\) is smooth
over \(S\) at those points, and the open subset \(U\) of points
of \(Y\) at which \(Y\) is smooth over \(S\) is an induced open
subgroup of \(Y\)
(cf.\ \hyperref[cor:sga3-VIB-2.3]{2.3}).
\begin{samepage}
The morphism
\[
  \tau:U\longrightarrow X
\]
is then an étale monomorphism by
\hyperref[prop:sga3-VIB-2.5]{2.5}, and hence an open immersion.
\end{samepage}
Since the fibers of \(X\) are connected and every open subgroup
of an algebraic group is also closed, \(\tau\) is a surjective
open immersion, that is, an isomorphism. Therefore \(U=X\), and
a fortiori \(Y=X\). This completes the proof of \hyperref[cor:sga3-VIB-6.5.4]{6.5.4}.

Proceeding as in
\hyperref[ex:sga3-VIB-6.2.4-e]{6.2.4\textup{(e)}}, one deduces
\hyperref[cor:sga3-VIB-6.5.5]{Corollary 6.5.5} from \hyperref[cor:sga3-VIB-6.5.4]{6.5.4}.
